% !TeX root = ../presentacion.tex
%=======================================================================
%  config/preambulo.tex — paquetes e idioma
%
%  Nada de este archivo es específico del proyecto. El aspecto visual vive
%  en config/tema.tex y los mecanismos de la plantilla en config/plantilla.tex.
%=======================================================================

%----------------------- Motor, idioma y tipografía --------------------
% Metropolis usa Fira Sans vía fontspec: hay que compilar con lualatex o
% xelatex. El Makefile ya lo hace. Con pdflatex la compilación aborta.
\usepackage{fontspec}
\usepackage[spanish,es-tabla,es-noquoting]{babel}
\usepackage{csquotes}

% Sin microtype a propósito: la versión de microtype que trae esta
% distribución usa \partokencontext, que el LuaTeX instalado no conoce, y la
% compilación se llena de «Undefined control sequence». En una lámina, con
% cuatro líneas de texto, el ajuste micrométrico no se echa de menos.

\usepackage{amsmath}
\usepackage{amssymb}

%------------------------------- Tablas --------------------------------
\usepackage{booktabs}
\usepackage{array}
\usepackage{tabularx}
\usepackage{ragged2e}
\usepackage{multirow}
\newcolumntype{L}[1]{>{\RaggedRight\arraybackslash}p{#1}}
\newcolumntype{C}[1]{>{\Centering\arraybackslash}p{#1}}
\newcolumntype{Y}{>{\RaggedRight\arraybackslash}X}
\renewcommand{\arraystretch}{1.2}

%--------------------------- Gráficos y color --------------------------
\usepackage{graphicx}
\usepackage{adjustbox}
% Los diagramas C4 no se duplican: se leen del informe (E1), que es donde
% se generan. Las capturas de pantalla propias de la charla van en figuras/.
\graphicspath{{figuras/}{../informe/figuras/}}

\usepackage{tcolorbox}
\tcbuselibrary{skins}

\usepackage{pifont}
\usepackage{fontawesome5}

% relsize da \smaller: «un paso menos que el tamaño actual». Es lo que hace
% que \id{...} encoja RESPECTO al texto que lo rodea. Con un \small fijo, un
% identificador dentro de una lista en \scriptsize salía más grande que la
% propia lista.
\usepackage{relsize}

%------------------------------ Listas ---------------------------------
\usepackage{enumitem}
\setlist{itemsep=4pt,parsep=0pt,topsep=4pt}

%------------------------------ Columnas -------------------------------
% Metropolis ya carga la maquinaria de columnas de beamer; esto solo fija
% un separador razonable para las láminas a dos columnas.
\setlength{\columnsep}{16pt}

%--------------------------- Código fuente -----------------------------
% Mismos colores que el informe, para que un fragmento se reconozca igual
% en el PDF y en la pantalla.
\usepackage{listings}

\definecolor{codeFondo}{HTML}{F7F7F7}
\definecolor{codeComentario}{HTML}{6A737D}
\definecolor{codePalabra}{HTML}{0B5394}
\definecolor{codeCadena}{HTML}{9C27B0}
\definecolor{codeBorde}{HTML}{DDDDDD}

\lstdefinestyle{estilobase}{
    backgroundcolor = \color{codeFondo},
    basicstyle      = \ttfamily\scriptsize,
    commentstyle    = \color{codeComentario}\itshape,
    keywordstyle    = \color{codePalabra}\bfseries,
    stringstyle     = \color{codeCadena},
    numbers         = none,          % en pantalla, los números estorban
    frame           = single,
    rulecolor       = \color{codeBorde},
    breaklines      = true,
    showstringspaces= false,
    tabsize         = 2,
    captionpos      = b,
    xleftmargin     = 6pt,
    literate =
      {á}{{\'a}}1 {é}{{\'e}}1 {í}{{\'i}}1 {ó}{{\'o}}1 {ú}{{\'u}}1
      {Á}{{\'A}}1 {É}{{\'E}}1 {Í}{{\'I}}1 {Ó}{{\'O}}1 {Ú}{{\'U}}1
      {ñ}{{\~n}}1 {Ñ}{{\~N}}1 {ü}{{\"u}}1 {¿}{{?`}}1 {¡}{{!`}}1
}
\lstset{style=estilobase}

\lstdefinelanguage{yaml}{
    keywords        = {true,false,null,version,services,image,build,ports,
                       environment,volumes,networks,depends_on,healthcheck,name},
    keywordstyle    = \color{codePalabra}\bfseries,
    sensitive       = false,
    comment         = [l]{\#},
    morestring      = [b]",
    morestring      = [b]'
}
% listings no trae JavaScript/TypeScript de fábrica: sin esta definición,
% un lstlisting con language=jsts aborta la compilación.
\lstdefinelanguage{jsts}{
    keywords        = {const,let,var,import,from,export,default,function,
                       return,true,false,null,undefined,async,await,new},
    keywordstyle    = \color{codePalabra}\bfseries,
    ndkeywords      = {name,remotes,shared,singleton},
    ndkeywordstyle  = \color{codeCadena}\bfseries,
    sensitive       = true,
    comment         = [l]{//},
    morecomment     = [s]{/*}{*/},
    morestring      = [b]",
    morestring      = [b]',
    morestring      = [b]`
}
\lstdefinelanguage{sql-pg}{
    language        = SQL,
    morekeywords    = {EXCLUDE,USING,gist,tstzrange,WITH,CONSTRAINT,DEFERRABLE},
    sensitive       = false
}

%----------------------- Numeración del apéndice -----------------------
% Sin esto, las láminas de respaldo inflan el «12 / 40» de la esquina y la
% barra de progreso: el público cree que faltan diez minutos más.
\usepackage{appendixnumberbeamer}
